\chapter{Quelltext der Use-Case-1-Umgebung}
\label{chap:quelltext-uc1}

Dieser Anhang enthält die vollständigen Dateien, mit denen die Sandbox von Use Case 1 definiert und
provisioniert wird. Die im Hauptteil (Abschnitt~\ref{sec:uc1-implementierung}) gezeigten Auszüge
sind hier vollständig abgedruckt, sodass sich die Umgebung nachbauen lässt. Lange Zeilen werden dabei
automatisch umgebrochen.

\section{Topologie (\texttt{topology.yml})}

\begin{GenericCode}
name: uc1-vulnscan
hosts:
  - name: target
    base_box:
      image: ubuntu-noble-x86_64
      man_user: ubuntu
    flavor: standard.small
  - name: scanner
    base_box:
      image: ubuntu-noble-x86_64
      man_user: ubuntu
    flavor: m1.medium
routers:
  - name: router
    base_box:
      image: debian-12-x86_64
      man_user: debian
    flavor: standard.small
networks:
  - name: lab-switch
    cidr: 192.168.20.0/24
    accessible_by_user: False
net_mappings:
  - host: target
    network: lab-switch
    ip: 192.168.20.5
  - host: scanner
    network: lab-switch
    ip: 192.168.20.10
router_mappings:
  - router: router
    network: lab-switch
    ip: 192.168.20.1
groups: []
\end{GenericCode}

\section{Abhängigkeiten (\texttt{requirements.yml})}

\begin{GenericCode}
roles:
  - name: user-access
    src: https://github.com/cyberrangecz/ansible-role-user-access.git
    scm: git
    version: v1.0.0
  - name: disable-qxl
    src: https://github.com/cyberrangecz/ansible-role-disable-qxl.git
    scm: git
    version: v1.0.0
  - name: sandbox-logging
    src: https://github.com/cyberrangecz/ansible-role-sandbox-logging.git
    scm: git
    version: v1.0.0

collections:
  - name: community.docker
  - name: ansible.posix
\end{GenericCode}

\section{Playbook (\texttt{playbook.yml})}

\begin{GenericCode}
---
- name: disable qxl
  hosts:
    - routers
    - hosts
  gather_facts: yes
  become: yes
  tasks:
    - include_role:
        name: disable-qxl
      when: ansible_os_family == 'Debian'

- name: set up vulnerable target
  hosts: target
  become: yes
  roles:
    - name: vulnerable-target

- name: set up scanner host
  hosts: scanner
  become: yes
  roles:
    - name: user-access
      user_access_username: user
      user_access_password: Password123
      user_access_sudo: True
    - name: scanner-host

- name: set up command logging
  hosts: hosts
  become: yes
  roles:
    - role: sandbox-logging
      slf_destination_port: '{% if hostvars["man"] is defined -%} 514 {%- else -%} 515 {%- endif %}'
...
\end{GenericCode}

\section{Rolle \texttt{vulnerable-target}}

\begin{GenericCode}
---
# Vulnerable Target für Use Case 1: Vulnerability Scanning
# Absichtliche Schwachstellen in verschiedenen Schwierigkeitsgraden
#
# Voraussetzung: Collection community.docker (in requirements.yml ergänzen):
#   collections:
#     - name: community.docker

# ============================================
# Basis-Dienste (Fehlkonfigurationen)
# ============================================

- name: Install vulnerable services
  apt:
    name:
      - vsftpd
      - telnetd
      - apache2
      - openssh-server
      - rsh-server
      - tftpd-hpa
      - samba
      - nfs-kernel-server
      - php
      - libapache2-mod-php
      - mysql-server
    state: present
    update_cache: yes

- name: Enable anonymous FTP access
  lineinfile:
    path: /etc/vsftpd.conf
    regexp: '^#?anonymous_enable'
    line: 'anonymous_enable=YES'

- name: Allow FTP write access
  lineinfile:
    path: /etc/vsftpd.conf
    regexp: '^#?write_enable'
    line: 'write_enable=YES'

- name: Create weak password users
  user:
    name: "{{ item.name }}"
    password: "{{ item.password | password_hash('sha512') }}"
    shell: /bin/bash
    groups: "{{ item.groups | default(omit) }}"
  loop:
    - { name: 'admin', password: 'admin' }
    - { name: 'test', password: 'test123' }
    - { name: 'backup', password: 'password' }
    - { name: 'guest', password: 'guest' }
    - { name: 'oracle', password: 'oracle' }

# ============================================
# Schwache SSH Konfiguration
# ============================================

- name: Weak SSH configuration
  blockinfile:
    path: /etc/ssh/sshd_config
    validate: '/usr/sbin/sshd -t -f %s'
    block: |
      # absichtlich schwach für UC1 - gültig für OpenSSH 9.x
      PermitRootLogin yes
      PasswordAuthentication yes
      PermitEmptyPasswords yes
      X11Forwarding yes
      MaxAuthTries 10
      Ciphers +aes128-cbc,3des-cbc
      MACs +hmac-sha1,hmac-md5
      KexAlgorithms +diffie-hellman-group1-sha1,diffie-hellman-group14-sha1
      HostKeyAlgorithms +ssh-rsa
      PubkeyAcceptedAlgorithms +ssh-rsa

# ============================================
# Apache Fehlkonfigurationen
# ============================================

- name: Insecure Apache configuration
  copy:
    dest: /etc/apache2/conf-available/insecure.conf
    content: |
      ServerTokens Full
      ServerSignature On
      <Directory /var/www/html>
        Options Indexes FollowSymLinks ExecCGI
        AllowOverride All
        Require all granted
      </Directory>
      TraceEnable On

- name: Activate insecure Apache config
  command: a2enconf insecure
  args:
    creates: /etc/apache2/conf-enabled/insecure.conf

# ============================================
# World-writable Dateien & SUID
# ============================================

- name: Create world-writable sensitive files
  file:
    path: "{{ item }}"
    state: touch
    mode: '0777'
  loop:
    - /etc/shadow.backup
    - /var/log/secret.log
    - /root/backup.tar.gz

- name: Set dangerous SUID bits
  file:
    path: "{{ item }}"
    mode: 'u+s'
  loop:
    - /bin/cp
    - /usr/bin/find
    - /usr/bin/vim.basic

# ============================================
# Sudo Fehlkonfigurationen
# ============================================

- name: Insecure sudoers configuration
  copy:
    dest: /etc/sudoers.d/insecure
    content: |
      test ALL=(ALL) NOPASSWD: /bin/bash
      admin ALL=(ALL) NOPASSWD: ALL
      backup ALL=(ALL) NOPASSWD: /usr/bin/vim, /usr/bin/find
    mode: '0440'
    validate: 'visudo -cf %s'

# ============================================
# Verwundbare Webanwendung
# ============================================

- name: Create vulnerable PHP application
  copy:
    dest: /var/www/html/admin.php
    content: |
      <?php
      // SQL Injection vulnerable
      $user = $_GET['user'];
      $query = "SELECT * FROM users WHERE name='$user'";
      // Command Injection vulnerable
      $cmd = $_GET['cmd'];
      if ($cmd) {
        echo shell_exec($cmd);
      }
      // Path Traversal vulnerable
      $file = $_GET['file'];
      if ($file) {
        include($file);
      }
      ?>

- name: Expose phpinfo
  copy:
    dest: /var/www/html/info.php
    content: |
      <?php phpinfo(); ?>

# ============================================
# Unsichere Samba Shares
# ============================================

- name: Insecure Samba configuration
  blockinfile:
    path: /etc/samba/smb.conf
    block: |
      [public]
        path = /tmp
        browseable = yes
        writable = yes
        guest ok = yes
        read only = no
        create mask = 0777
        directory mask = 0777
      [admin]
        path = /etc
        browseable = yes
        writable = yes
        guest ok = yes

# ============================================
# Unsichere NFS Exports
# ============================================

- name: Insecure NFS exports
  copy:
    dest: /etc/exports
    content: |
      /home *(rw,sync,no_root_squash,no_subtree_check)
      /tmp *(rw,sync,no_root_squash,insecure)
      /etc *(ro,sync,no_root_squash)

# ============================================
# MySQL Schwachstellen
# Hinweis: Diese Syntax (PASSWORD()/inline GRANT) ist unter MySQL 8.0
# nicht mehr gültig und schlägt fehl. Für einen funktionierenden
# DB-Stand auf CREATE USER + ALTER USER + GRANT umstellen.
# ============================================

- name: MySQL weak root configuration
  shell: |
    mysql -e "UPDATE mysql.user SET authentication_string=PASSWORD('root') WHERE User='root';"
    mysql -e "GRANT ALL PRIVILEGES ON *.* TO 'admin'@'%' IDENTIFIED BY 'admin' WITH GRANT OPTION;"
    mysql -e "FLUSH PRIVILEGES;"
  ignore_errors: yes

# ============================================
# Swap (2 GB) - nötig für die JVM-Container (Tomcat/Jenkins)
# auf der ressourcenschwachen Sandbox-VM
# ============================================

- name: Swap-File anlegen
  ansible.builtin.command: "fallocate -l 2G /swapfile"
  args:
    creates: /swapfile

- name: Rechte auf Swap-File setzen
  ansible.builtin.file:
    path: /swapfile
    owner: root
    group: root
    mode: '0600'

- name: Swap formatieren
  ansible.builtin.command: mkswap /swapfile
  register: _mkswap
  changed_when: "'Setting up swapspace' in _mkswap.stdout"
  failed_when:
    - _mkswap.rc != 0
    - "'already' not in _mkswap.stderr"

- name: Swap in /etc/fstab persistieren (überlebt Reboot)
  ansible.posix.mount:
    path: none
    src: /swapfile
    fstype: swap
    opts: sw
    state: present

- name: Swap aktivieren
  ansible.builtin.command: swapon /swapfile
  register: _swapon
  changed_when: _swapon.rc == 0
  failed_when:
    - _swapon.rc != 0
    - "'already' not in _swapon.stderr"

# ============================================
# Docker (für die CVE-Container)
# ============================================

- name: Docker installieren
  ansible.builtin.apt:
    name: docker.io
    state: present
    update_cache: true

- name: Docker-Dienst starten und aktivieren
  ansible.builtin.service:
    name: docker
    state: started
    enabled: true

- name: Python-Docker-SDK (für docker_container-Modul)
  ansible.builtin.apt:
    name: python3-docker
    state: present

# ============================================
# Verwundbare CVE-Dienste als Container
# A: Apache 2.4.49 (CVE-2021-41773), Tomcat 9.0.30 (CVE-2020-1938),
#    Grafana 8.3.0 (CVE-2021-43798), Jenkins 2.138.1 (CVE-2018-1000861)
# ============================================

- name: Verwundbare CVE-Container starten
  community.docker.docker_container:
    name: "{{ item.name }}"
    image: "{{ item.image }}"
    ports: "{{ item.ports }}"
    restart_policy: unless-stopped
    state: started
    pull: true
  loop:
    - { name: vuln-apache,  image: "httpd:2.4.49",            ports: ["8080:80"] }
    - { name: vuln-tomcat,  image: "tomcat:9.0.30",           ports: ["8090:8080", "8009:8009"] }
    - { name: vuln-grafana, image: "grafana/grafana:8.3.0",   ports: ["3000:3000"] }
    - { name: vuln-jenkins, image: "jenkins/jenkins:2.138.1", ports: ["8085:8080"] }

# ============================================
# Services aktivieren
# ============================================

- name: Start all vulnerable services
  service:
    name: "{{ item }}"
    state: started
    enabled: yes
  loop:
    - vsftpd
    - apache2
    - smbd
    - nmbd
    - nfs-kernel-server
    - mysql
  ignore_errors: yes

- name: Restart SSH
  service:
    name: ssh
    state: restarted
\end{GenericCode}

\section{Rolle \texttt{scanner-host}}

\begin{GenericCode}
---
# Scanner Host für Use Case 1
# OpenVAS via Greenbone Community Containers (apt-Pakete gibt es unter Noble nicht mehr).
# Image-/Feed-Download läuft non-blocking im Hintergrund, damit die Allocation
# nicht in einen Timeout läuft. Docker restart_policy hält den Stack über Reboots.

- name: Update apt cache
  apt:
    update_cache: yes
    cache_valid_time: 3600

- name: Install base tools
  apt:
    name:
      - curl
      - wget
      - nmap
      - net-tools
      - dnsutils
      - vim
      - htop
      - tmux
      - unzip
      - ca-certificates
      - gnupg
    state: present

# --- Docker Engine + Compose-Plugin (offizielles Docker-Repo) ---
- name: Create Docker keyring directory
  file:
    path: /etc/apt/keyrings
    state: directory
    mode: '0755'

- name: Add Docker GPG key
  get_url:
    url: https://download.docker.com/linux/ubuntu/gpg
    dest: /etc/apt/keyrings/docker.asc
    mode: '0644'

- name: Add Docker apt repository
  apt_repository:
    repo: "deb [arch=amd64 signed-by=/etc/apt/keyrings/docker.asc] https://download.docker.com/linux/ubuntu noble stable"
    filename: docker

- name: Install Docker
  apt:
    name:
      - docker-ce
      - docker-ce-cli
      - containerd.io
      - docker-compose-plugin
    state: present
    update_cache: yes

- name: Enable and start Docker
  systemd:
    name: docker
    enabled: yes
    state: started

- name: Add user to docker group
  user:
    name: ubuntu
    groups: docker
    append: yes

# --- Greenbone Community Edition (OpenVAS) ---
- name: Create Greenbone directory
  file:
    path: /opt/greenbone
    state: directory
    mode: '0755'

- name: Download Greenbone compose file
  get_url:
    url: https://greenbone.github.io/docs/latest/_static/compose.yaml
    dest: /opt/greenbone/compose.yaml
    mode: '0644'

- name: Pull and start Greenbone (non-blocking background job)
  shell: |
    docker compose -f /opt/greenbone/compose.yaml pull
    docker compose -f /opt/greenbone/compose.yaml up -d
  args:
    chdir: /opt/greenbone
  async: 3600
  poll: 0

# --- Nuclei (template-basierter Scanner) ---
- name: Download Nuclei
  get_url:
    url: https://github.com/projectdiscovery/nuclei/releases/download/v3.3.5/nuclei_3.3.5_linux_amd64.zip
    dest: /tmp/nuclei.zip
    mode: '0644'

- name: Extract Nuclei
  unarchive:
    src: /tmp/nuclei.zip
    dest: /usr/local/bin/
    remote_src: yes
    creates: /usr/local/bin/nuclei

- name: Make Nuclei executable
  file:
    path: /usr/local/bin/nuclei
    mode: '0755'
    owner: root
    group: root

# --- Arbeitsverzeichnis + README ---
- name: Create scanner working directory
  file:
    path: /home/user/scans
    state: directory
    owner: user
    group: user
    mode: '0755'

- name: Create README for the scanner host
  copy:
    dest: /home/user/README.md
    owner: user
    group: user
    content: |
      # UC1 Scanner Host

      ## OpenVAS / Greenbone (Container)
      OpenVAS laeuft als Greenbone-Community-Container-Stack.
      Nach dem ersten Hochfahren lädt der Scanner die Feeds im Hintergrund
      (CVE/SCAP/VTs) - das dauert beim ersten Mal ca. 15-30 Minuten.

      Status prüfen:
        docker compose -f /opt/greenbone/compose.yaml ps
        docker compose -f /opt/greenbone/compose.yaml logs -f gvmd

      Admin-Passwort setzen (Standard admin/admin):
        docker compose -f /opt/greenbone/compose.yaml exec -u gvmd gvmd \
          gvmd --user=admin --new-password='DEIN_PASSWORT'

      Web-Interface (bindet nur auf 127.0.0.1):
        SSH-Tunnel vom Laptop, dann https://localhost:9392
          ssh -L 9392:127.0.0.1:9392 user@<scanner>

      ## Nuclei
        nuclei -u http://192.168.20.5
        nuclei -target 192.168.20.5 -severity critical,high

      ## Nessus Essentials (manuell)
      1. Account: https://www.tenable.com/products/nessus/nessus-essentials
      2. .deb: https://www.tenable.com/downloads/nessus
      3. sudo dpkg -i Nessus-*.deb && sudo systemctl start nessusd
      4. Browser: https://localhost:8834

      ## Ziel-System
      Target: 192.168.20.5

      ## Scans speichern in
      /home/user/scans/
\end{GenericCode}
